\chapter{Inverse Kinematics}
\label{ch:ik}

% =====================================================================
%  Inverse kinematics: 2R planar
%  Goes into chapters/04_inverse_kinematics.tex
%  Requires: tikz (calc, angles, quotes), subcaption, siunitx, cleveref
%  References \cref{lem:triangle-laws} from the prerequisites.
% =====================================================================

\section{Inverse kinematics: 2R}
\label{sec:ik-2r-planar}

\begin{figure}[tbp]
  \centering
  % Illustrative geometry shared by both sketches
  \pgfmathsetmacro{\Ra}{2}       % R_1
  \pgfmathsetmacro{\Rb}{1.6}     % R_2
  \pgfmathsetmacro{\dT}{3}       % distance base--target
  \pgfmathsetmacro{\th}{45}      % theta
  \pgfmathsetmacro{\al}{acos((\Ra^2+\dT^2-\Rb^2)/(2*\Ra*\dT))} % alpha
  %
  \begin{subfigure}[b]{0.48\textwidth}
    \centering
    \begin{tikzpicture}[scale=1.4, thick]
      \coordinate (O)  at (0,0);
      \coordinate (Y)  at (1,0);
      \coordinate (Z0) at (0,1);
      \coordinate (T)  at (\th:\dT);
      \coordinate (J)  at (\th+\al:\Ra);
      \coordinate (X)  at ($(O)!1.45!(J)$);    % extension of link 1
      % ground, helper lines
      \draw (-0.8,0) -- (3,0);
      \draw[dashed, thin] (O) -- (0,2.4);
      \draw[dashed, thin] (J) -- (X);
      \draw[dashed, thin] (O) -- (T);
      % links
      \draw (O) -- node[left=3pt, pos=0.35] {$R_1$} (J);
      \draw (J) -- node[above, pos=0.7] {$R_2$} (T);
      \fill (J) circle (1.5pt);
      \node at (T) {$\times$};
      \node[right] at (T) {$\vect{t}$};
      % angles at the base
      \pic[draw, thin, "$\theta$", angle radius=14pt, angle eccentricity=1.45]
        {angle = Y--O--T};
      \pic[draw, thin, "$\alpha$", angle radius=26pt, angle eccentricity=1.3]
        {angle = T--O--J};
      \pic[draw, thin, "$\varphi_1$", angle radius=50pt, angle eccentricity=1.2]
        {angle = J--O--Z0};
      % angles at the elbow and the target
      \pic[draw, thin, "$\beta$", angle radius=10pt, angle eccentricity=1.8]
        {angle = O--J--T};
      \pic[draw, thin, "$\varphi_2$", angle radius=14pt, angle eccentricity=1.6]
        {angle = T--J--X};
      \pic[draw, thin, "$\gamma$", angle radius=14pt, angle eccentricity=1.6]
        {angle = J--T--O};
    \end{tikzpicture}
    \caption{Configuration 1 with all angles.}
    \label{fig:ik-2r-angles}
  \end{subfigure}\hfill
  \begin{subfigure}[b]{0.48\textwidth}
    \centering
    \begin{tikzpicture}[scale=1.15, thick]
      \coordinate (O)  at (0,0);
      \coordinate (T)  at (\th:\dT);
      \coordinate (J1) at (\th+\al:\Ra);
      \coordinate (J2) at (\th-\al:\Ra);
      % circles whose intersections are the possible elbow positions
      \draw[thin, gray] (\Ra,0) arc[start angle=0, end angle=180, radius=\Ra];
      \draw[thin, gray] (T) circle[radius=\Rb];
      \draw (-2.2,0) -- (3.9,0);
      \draw[dashed, thin] (O) -- (T);
      % configuration 1 (solid) and 2 (dashed)
      \draw (O) -- (J1) -- (T);
      \draw[densely dashed] (O) -- (J2) -- (T);
      \fill (J1) circle (1.5pt) node[above left]  {1};
      \fill (J2) circle (1.5pt) node[below right] {2};
      \node at (T) {$\times$};
      \node[above right] at (T) {$\vect{t}$};
    \end{tikzpicture}
    \caption{Both configurations reaching $\vect{t}$.}
    \label{fig:ik-2r-configs}
  \end{subfigure}
  \caption{Inverse kinematics of the planar 2R arm. The possible elbow
    positions are the intersections of the circle of radius $R_1$ about
    the base with the circle of radius $R_2$ about the target.}
  \label{fig:ik-2r}
\end{figure}

\enquote{Here the second configuration does not exist}

% Added during transcription -- please check.
As shown in \cref{fig:ik-2r-angles}, $\alpha$, $\beta$ and $\gamma$ are the
interior angles at the base, the elbow and the target of the triangle
formed by the two links and the line from the base to the target, and
$\theta$ is the angle between the ground and that line.

$\alpha$, $\beta$ and $\gamma$ can be determined via the laws of sine and
cosine. Of course, we only need $\alpha$ and $\beta$:
\begin{align*}
  \alpha &= \arctan\left(\frac{\sin(\alpha)}{\cos(\alpha)}\right), \\
  \beta  &= \arctan\left(\frac{\sin(\beta)}{\cos(\beta)}\right),
\end{align*}
where we find $\sin(\alpha)$ and $\cos(\alpha)$ via
\cref{lem:triangle-laws}. $\theta$ we can find directly via
\[
  \theta = \arctan\left(\frac{t_2}{t_1}\right)
\]
with the target $\vect{t} = (t_1, t_2)\T$.

Now we have all the ingredients we need to find $\varphi_1$ and
$\varphi_2$ for both configurations:
\begin{enumerate}
  \item $\varphi_1 = \theta + \alpha - \ang{90}$ \quad and \quad
        $\varphi_2 = -\ang{180} + \beta$,
  \item $\varphi_1 = \theta - \alpha - \ang{90}$ \quad and \quad
        $\varphi_2 = \ang{180} - \beta$.
\end{enumerate}

\paragraph{Hint.}
Most libraries offer an \texttt{atan2}, which resolves all four quadrants.

\paragraph{Reachability.}
We assume $R_1 > R_2$. Then a target $\vect{t} = (t_1, t_2)\T$ is only
reachable if
\[
  R_1 - R_2 \le \lvert \vect{t} \rvert \le R_1 + R_2 ,
\]
where the equals disappears in practice: on the boundary the arm is
either fully stretched or fully folded, both configurations coincide,
and small errors in the link lengths or rounding in the computation
already push the target out of reach.

% =====================================================================
%  chapters/04_inverse_kinematics.tex -- 4R arm and 5R extension
%  Goes after the section "Inverse kinematics: 2R planar".
%  Requires: tikz (calc, angles, quotes, arrows.meta), subcaption
% =====================================================================

\section{Inverse kinematics: 4R arm}
\label{sec:4r-ik}

The arm and its variables are introduced in \cref{sec:4r-fk}.

\paragraph{Base angle $\varphi_0$.}
The angle of the $x$-$y$-position of the target
$\vect{t} = (1, t_1, t_2, t_3)\T$ we can compute directly:
\[
  \varphi_0' = \arctan\left(\frac{t_2}{t_1}\right) .
\]
But we need to correct this angle by $\delta\varphi_0$, the angle between
the arm's $x'$-axis and the direction to the target (see
\cref{fig:4r-structure}). To compute $\delta\varphi_0$ we need the
$x$-$y$-length of the arm:
\[
  \delta\varphi_0 = \arcsin\left(\frac{D}{\sqrt{t_1^2 + t_2^2}}\right)
  \qquad\text{with } D = D_0 + D_1 + D_2 .
\]
$\delta\varphi_0 \in [0, \frac{\pi}{2})$, so $\arcsin$ already gives the
unique answer. It also introduces an additional reachability condition:
\[
  D \le \sqrt{t_1^2 + t_2^2} .
\]

\paragraph{Two configurations for the base angle.}
As long as the joint angles allow it, one can always create a
configuration in which the $D_i$ all point in the positive direction and
one in which they all point in the negative direction
(\cref{fig:4r-base-configs}). For a given target this means that two
configurations are possible. One of them requires the base angle
\begin{equation}
  \boxed{\varphi_0^{(1)} = \varphi_0' - \delta\varphi_0} ,
  \label{eq:4r-phi0-1}
\end{equation}
where $\varphi_0'$ is the angle of the target and $\delta\varphi_0$
describes the angle of the arm configuration. The other first requires a
rotation by $\pi$, followed by an inversion of all planar joint angles,
$\varphi_i \to -\varphi_i$ for $i = 1, 2, 3$. Afterwards, as \cref{fig:4r-base-configs} shows, we have to rotate by
$\delta\varphi_0$ further: $\varphi_0^{(2)} = \pi + \varphi_0' +
\delta\varphi_0$. Since this number is greater than $\pi$,
\begin{equation}
  \boxed{\varphi_0^{(2)} = \varphi_0' - \pi + \delta\varphi_0} .
  \label{eq:4r-phi0-2}
\end{equation}

\begin{figure}[tbp]
  \centering
  % Same target t = (3, 0), D = 1.2 (exaggerated); delta phi_0 = asin(0.4).
  \begin{tikzpicture}[thick, scale=1.5]
    \pgfmathsetmacro{\dl}{asin(0.4)}
    \draw[-{Stealth[length=4pt]}, semithick] (-1.4,1.2) -- (-0.9,1.2) node[right] {$x$};
    \draw[-{Stealth[length=4pt]}, semithick] (-1.4,1.2) -- (-1.4,1.7) node[above] {$y$};
    \draw[thin] (0,0) circle (0.3);
    \fill (0,0) circle (1.3pt);
    \draw[dotted, thin] (0,0) -- (3,0);
    % configuration 1: phi_0 = -delta, offsets along +y'
    \begin{scope}[rotate=-\dl]
      \draw (0,0) -- (0,0.4) -- (1,0.4) -- (1,0.8) -- (2,0.8) -- (2,1.2) -- (2.75,1.2);
      \draw[-{Stealth[length=4pt]}, gray, semithick] (0,0) -- (0.9,0)
        node[pos=1, below, black] {$x'^{(1)}$};
      \node[above, font=\small] at (1.5,0.8) {1};
    \end{scope}
    % configuration 2: phi_0 = -pi + delta, links reach along -x'
    \begin{scope}[rotate=180+\dl]
      \draw[densely dashed] (0,0) -- (0,0.4) -- (-1,0.4) -- (-1,0.8) -- (-2,0.8)
        -- (-2,1.2) -- (-2.75,1.2);
      \draw[-{Stealth[length=4pt]}, gray, semithick] (0,0) -- (0.9,0)
        node[pos=1, below, black] {$x'^{(2)}$};
      \node[below, font=\small] at (-1.5,0.8) {2};
    \end{scope}
    \node at (3,0) {$\times$};
    \node[right] at (3,0) {$\vect{t}$};
  \end{tikzpicture}
  \caption{Top view of the two base configurations reaching the same
    target: configuration 1 (solid) and configuration 2 (dashed), whose
    $x'$-axis points away from the links.}
  \label{fig:4r-base-configs}
\end{figure}

\paragraph{Reduction to a plane.}
Once we have $\varphi_0$, we can calculate the remaining angles by only
considering a plane, specifically the $z$-$x'$-plane, where $x'$ is the
arm's own horizontal axis shown in \cref{fig:4r-base-configs}. Why is
that the axis? The position in the $y'$-direction cannot be affected,
since the $D_i$ cannot be changed. In the $x'$-direction, we get changes
by moving the links according to the joints' degrees of freedom.
The target position in the $x'$-$z$-plane is
\[
  \vect{t}'^{(1)} =
  \begin{pmatrix} \sqrt{t_1^2 + t_2^2 - D^2} \\ t_3 \end{pmatrix} .
\]
In the case of the second configuration, the $x'$-axis points in the
opposite direction of the links. Therefore we have to invert the sign of
the $x'$-component, where $\varphi_0^{(2)}$ is not explicitly included in
the formula:
\[
  \vect{t}'^{(2)} =
  \begin{pmatrix} -\sqrt{t_1^2 + t_2^2 - D^2} \\ t_3 \end{pmatrix} .
\]

\paragraph{Wrist position in the $x'$-$z$-plane.}
The 3R problem in the $x'^{(i)}$-$z$-plane looks identical for both
configurations, with the only difference being that the $x'^{(2)}$-axis is
mirrored on the $z$-axis (\cref{fig:4r-planar}).

\begin{figure}[tbp]
  \centering
  \begin{subfigure}[b]{0.56\textwidth}
    \centering
    \begin{tikzpicture}[thick, scale=1.2]
      \draw (-1.4,0) -- (3.6,0);
      \draw (-0.4,0) rectangle (0.4,0.4);
      \draw[|<->|, thin] (0.6,0) -- node[right] {$h$} (0.6,0.4);
      % bases of both configurations
      \draw[-{Stealth[length=4pt]}, semithick, blue!70!black]
        (0,0) -- (1.1,0) node[below] {$x'^{(1)}$};
      \draw[-{Stealth[length=4pt]}, semithick, blue!70!black]
        (0,0) -- (-1.1,0) node[below] {$x'^{(2)}$};
      \draw[-{Stealth[length=4pt]}, semithick, blue!70!black]
        (0,0) -- (0,1.0) node[left] {$z$};
      \coordinate (S) at (0,0.4);
      \coordinate (E) at ($(S)+(70:1.6)$);
      \coordinate (W) at ($(E)+(-5:1.6)$);
      \coordinate (T) at ($(W)+(-55:0.9)$);
      \draw[dashed, thin] (S) -- node[below right, pos=0.6] {$\lVert\vect{p}_3''\rVert$} (W);
      \draw (S) -- node[left, pos=0.8] {$R_1$} (E)
                -- node[above] {$R_2$} (W)
                -- node[right] {$R_3$} (T);
      \foreach \p in {S,E,W} \fill (\p) circle (1.3pt);
      \node at (T) {$\times$};
    \end{tikzpicture}
    \caption{The planar problem with the bases of both configurations.}
    \label{fig:4r-planar}
  \end{subfigure}\hfill
  \begin{subfigure}[b]{0.4\textwidth}
    \centering
    \begin{tikzpicture}[thick, scale=1.2]
      \coordinate (O) at (0,0);
      \coordinate (Z) at (0,1.6);
      % v and its rotation, in basis 1 (right) and mirrored in basis 2 (left)
      \coordinate (V1) at (72:1.2);
      \coordinate (W1) at (32:1.2);
      \coordinate (V2) at (108:1.2);
      \coordinate (W2) at (148:1.2);
      \draw[thin, gray] (1.2,0) arc[start angle=0, end angle=180, radius=1.2];
      \draw[-{Stealth[length=4pt]}, semithick, blue!70!black]
        (O) -- (1.6,0) node[below] {$x'^{(1)}$};
      \draw[-{Stealth[length=4pt]}, semithick, blue!70!black]
        (O) -- (-1.6,0) node[below] {$x'^{(2)}$};
      \draw[-{Stealth[length=4pt]}, semithick, blue!70!black]
        (O) -- (Z) node[left] {$z$};
      \draw[-{Stealth[length=4pt]}] (O) -- (V1) node[above right=-2pt] {$\vect{v}^{(1)}$};
      \draw[-{Stealth[length=4pt]}] (O) -- (W1) node[right] {$\tilde{\vect{v}}^{(1)}$};
      \draw[-{Stealth[length=4pt]}] (O) -- (V2) node[above left=-2pt] {$\vect{v}^{(2)}$};
      \draw[-{Stealth[length=4pt]}] (O) -- (W2) node[left] {$\tilde{\vect{v}}^{(2)}$};
      \pic[draw, thin, {Stealth[length=3pt]}-, "$\lvert\varphi\rvert$",
           angle radius=22pt, angle eccentricity=1.55] {angle = W1--O--V1};
      \pic[draw, thin, -{Stealth[length=3pt]}, "$\lvert\varphi\rvert$",
           angle radius=22pt, angle eccentricity=1.55] {angle = V2--O--W2};
    \end{tikzpicture}
    \caption{A vector $\vect{v}$ and its rotation $\tilde{\vect{v}}$,
      described in the basis of configuration 1 and in the mirrored basis
      of configuration 2.}
    \label{fig:4r-handedness}
  \end{subfigure}
  \caption{The arm plane of both configurations.}
\end{figure}

Obviously, in both bases the magnitude of the angles does not change;
however, the \enquote{handedness} of the rotation changes, as
\cref{fig:4r-handedness} visualizes. Therefore we only need to compute
$\varphi_1$, $\varphi_2$ and $\varphi_3$ for one case and find the angles
for the other case by inverting the handedness:
\begin{equation}
  \varphi_i^{(2)} = -\varphi_i^{(1)} , \qquad i = 1, 2, 3 .
  \label{eq:4r-config2}
\end{equation}

The target comes with an additional constraint, the angle of vertical
approach $\theta$, where $\theta = 0$ implies approaching from straight
above (\cref{fig:4r-approach}). With the angle of approach, one can find
the position $\vect{p}_3'$ of the third joint:
\[
  \vect{p}_3'^{(1)} = \vect{t}'^{(1)}
    + R_3 \begin{pmatrix} -\sin\theta \\ \cos\theta \end{pmatrix} .
\]

\begin{figure}[tbp]
  \centering
  \begin{tikzpicture}[thick, scale=1.3]
    \draw[-{Stealth[length=4pt]}, semithick] (-1.3,1.2) -- (-0.8,1.2) node[right] {$x'$};
    \draw[-{Stealth[length=4pt]}, semithick] (-1.3,1.2) -- (-1.3,1.7) node[above] {$z$};
    \coordinate (T) at (0,0);
    \coordinate (W) at (125:1.1);
    \coordinate (U) at (0,1.3);
    \draw[thin, gray] (T) circle (1.1);
    \draw[dashed, thin] (T) -- (U);
    \draw (T) -- node[below left] {$R_3$} (W);
    \fill (W) circle (1.3pt);
    \node at (T) {$\times$};
    \pic[draw, thin, "$\theta$", angle radius=18pt, angle eccentricity=1.4]
      {angle = U--T--W};
  \end{tikzpicture}
  \caption{Angle of approach $\theta$.}
  \label{fig:4r-approach}
\end{figure}

\paragraph{Angles $\varphi_1$ and $\varphi_2$.}
Now we can calculate the angles $\varphi_1$ and $\varphi_2$ the same way
as for the 2R arm (\cref{sec:ik-2r-planar}), with $\vect{p}_3''$
replacing $\vect{t}$ (\cref{fig:4r-planar}), where
\[
  \vect{p}_3'' = \vect{p}_3' - \begin{pmatrix} 0 \\ h \end{pmatrix}
\]
is the wrist position relative to the shoulder.

\paragraph{Angle $\varphi_3$.}
With $\Phi = \varphi_1 + \varphi_2 + \varphi_3$ the absolute angle of
the last link (\cref{fig:4r-phi3}; the angles in the drawing are
unsigned), we find for the first configuration
\[
  \Phi - \theta = \pi ,
  \qquad
  \varphi_3^{(1)} = \pi + \theta - \varphi_1^{(1)} - \varphi_2^{(1)} .
\]
For $\varphi_3 > \pi$, interpret $\varphi_3 \to \varphi_3 - 2\pi$.
% Added during transcription -- follows from eq:4r-config2, please check.
For the second configuration, \eqref{eq:4r-config2} gives
\[
  \varphi_3^{(2)} = -\pi - \theta - \varphi_1^{(2)} - \varphi_2^{(2)} ,
\]
which accordingly has to be wrapped into $(-\pi, \pi]$ by adding $2\pi$
for $\varphi_3 \le -\pi$.

\begin{figure}[tbp]
  \centering
  \begin{tikzpicture}[thick, scale=1.2]
    \coordinate (W) at (0,0);
    \coordinate (P) at (150:1.6);           % previous link (to elbow)
    \coordinate (T) at (-60:1.2);           % last link, down and outwards
    \coordinate (U) at (0,1);
    \coordinate (Dn) at (0,-1.4);
    \draw[dashed, thin] (U) -- (Dn);
    \draw (P) -- (W) -- (T);
    \fill (W) circle (1.3pt);
    \node at (T) {$\times$};
    \pic[draw, thin, "$\Phi$", angle radius=14pt, angle eccentricity=1.7]
      {angle = T--W--U};
    \pic[draw, thin, "$\theta$", angle radius=26pt, angle eccentricity=1.3]
      {angle = Dn--W--T};
  \end{tikzpicture}
  \caption{Absolute angle $\Phi$ of the last link, measured from the
    upward vertical, and the angle of approach $\theta$. Both are drawn
    unsigned; their signs follow from the handedness of the respective
    basis.}
  \label{fig:4r-phi3}
\end{figure}

Combining the two base configurations with the two solutions of the
planar problem (elbow up and elbow down), the 4R arm, in the version used
here, can choose between up to four configurations for a given set of
constraints, as far as the joint limits allow it.

\subsection{Extension: 5R -- LeRobot style}
\label{sec:5r}

LeRobot style is basically the 4R arm that was just introduced, but
additionally the end effector, a gripper, can rotate about the last
link's axis. The end effector of LeRobot is $\pi$-symmetric.

The local basis $\vect{g}_x$, $\vect{g}_y$, $\vect{g}_z$ at the gripper
is described by $\vect{g}_z$ aligning with the last link's axis,
$\vect{g}_y$ being the tangent to the base circle at angle $\varphi_0$,
and $\vect{g}_x$ being perpendicular to both, with a sign such that for
$\theta = 0$, $\vect{g}_x$ points in radial direction outwards:
\[
  \vect{g}_x = \begin{pmatrix} \cos\theta\cos\varphi_0 \\
                               \cos\theta\sin\varphi_0 \\
                               \sin\theta \end{pmatrix} , \qquad
  \vect{g}_y = \begin{pmatrix} -\sin\varphi_0 \\ \cos\varphi_0 \\ 0 \end{pmatrix} , \qquad
  \vect{g}_z = \begin{pmatrix} -\sin\theta\cos\varphi_0 \\
                               -\sin\theta\sin\varphi_0 \\
                               \cos\theta \end{pmatrix} .
\]
For the second configuration, the same formulas if one
ignores the initial rotation by $\pi$, i.e.\ with $\varphi_0 = \varphi_0^{(2)} - \pi$
and $\theta$ unchanged. This however does also neglect that 
we turn the arm upside down after the initial rotation by $\pi$, 
so the angle $\chi$ computed below has to be
turned back: $\chi^{(2)} = \chi - \pi$. For the $\pi$-symmetric gripper
of LeRobot, this correction can be omitted.

Now we can find the line along which the grabbing motion is supposed to
take place in the new basis. With $\vect{l}$ the direction of the
grabbing motion, we look for $a$ and $b$ with
\[
  a \, \vect{g}_x + b \, \vect{g}_y = \vect{l} .
\]
This system is overdetermined, and the input is likely noisy; therefore
solving it exactly is most likely not possible. Since $\vect{g}_x$ and
$\vect{g}_y$ are orthonormal, we use
\[
  a = \vect{g}_x\T \vect{l} , \qquad
  b = \vect{g}_y\T \vect{l} ,
\]
which is the least-squares solution, i.e.\ the exact solution for the
projection of $\vect{l}$ onto the $\vect{g}_x$-$\vect{g}_y$-plane. Now we
can calculate the angle in the $\vect{g}_x$-$\vect{g}_y$-plane:
\[
  \chi = \arctan\left(\frac{b}{a}\right) .
\]
This is the angle by which we have to rotate the end effector about the
last link's axis.

\paragraph{Example: grabbing from above along the $x$-axis.}
With $\theta = 0$ and $\vect{l} = (1, 0, 0)\T$:
\[
  \vect{g}_x = \begin{pmatrix} \cos\varphi_0 \\ \sin\varphi_0 \\ 0 \end{pmatrix} ,
  \qquad
  \vect{g}_y = \begin{pmatrix} -\sin\varphi_0 \\ \cos\varphi_0 \\ 0 \end{pmatrix} ,
  \qquad
  \begin{pmatrix} a \\ b \end{pmatrix}
  = \begin{pmatrix} \cos\varphi_0 \\ -\sin\varphi_0 \end{pmatrix} ,
\]
\[
  \chi = \arctan\left(\frac{-\sin\varphi_0}{\cos\varphi_0}\right) = -\varphi_0 .
\]